pdfoutput=1
\pdfoutput=1
% !TEX program = xelatex
% ===========================================================================
%  SCX Environmental Intergenerational Potential:
%  Climate Change as Cross-Generational Potential Jump —
%  Theorem 10 (Confinement) at Civilization Scale
%  Date: 2026-07-01
% ===========================================================================
\documentclass[12pt,a4paper]{article}
% ---- Mathematics ----
\usepackage{amsmath,amssymb,amsthm,mathtools}
\usepackage{bm}
\usepackage{bbm}
\usepackage{dsfont}
% ---- Theorem Environments ----
\newtheorem{theorem}Theorem[section]
\newtheorem{proposition}[theorem]{Proposition}
\newtheorem{corollary}[theorem]Corollary
\newtheorem{lemma}[theorem]Lemma
\newtheorem{definition}[theorem]Definition
\newtheorem{assumption}[theorem]Assumption
\theoremstyle{remark}
\newtheorem{remark}[theorem]Remark
\newtheorem{example}[theorem]{Example}
% ---- Layout ----
\usepackage[margin=2.5cm]{geometry}
\usepackage{booktabs}
\usepackage{graphicx}
\usepackage{hyperref}
\usepackage{enumitem}
\usepackage[capitalise,noabbrev]{cleveref}
\usepackage{xcolor}
\usepackage{tikz}
\usepackage{float}
\usepackage{algorithm}
\usepackage{algpseudocode}
% ---- Hyperref ----
\hypersetup{
    colorlinks=true,
    citecolor=blue,
    urlcolor=blue,
    pdftitle={SCX Environmental Intergenerational Potential — Climate Change as Intergenerational Potential Jump},
    pdfauthor={SCX},
    pdfsubject={SCX Environmental Gauge Theory},
}
% ---- Honest critique markers ----
\newtheorem{honestattack}{Honest Strike}
\renewcommand{\thehonestattack}{\textcolor{red}{\textbf{[!] Strike \arabic{honestattack}}}}
\newenvironment{attackbox}{%
  \begin{center}\color{red}\begin{minipage}{0.92\textwidth}\begin{honestattack}}{%
  \end{honestattack}\end{minipage}\end{center}}
% ---- Math Operators ----
\DeclareMathOperator{\PE}{PE}
\DeclareMathOperator{\Cov}{Cov}
\DeclareMathOperator{\Var}{Var}
\DeclareMathOperator{\KL}{D_{KL}}
\DeclareMathOperator{\TV}{TV}
\DeclareMathOperator*{\argmin}{arg\,min}
\DeclareMathOperator*{\argmax}{arg\,max}
\DeclareMathOperator{\sgn}{sgn}
\DeclareMathOperator{\Tr}{Tr}
\DeclareMathOperator{\diag}{diag}
\DeclareMathOperator{\spec}{spec}
% ---- Custom notations ----
\newcommand{\R}{\mathbb{R}}
\newcommand{\N}{\mathbb{N}}
\newcommand{\C}{\mathbb{C}}
\newcommand{\E}{\mathbb{E}}
\newcommand{\Pbb}{\mathbb{P}}
\newcommand{\Y}{\mathcal{Y}}
\newcommand{\X}{\mathcal{X}}
\newcommand{\Sstates}{\mathcal{S}}
\newcommand{\F}{\mathcal{F}}
\newcommand{\D}{\mathcal{D}}
\newcommand{\Gauge}{\mathcal{G}}
\newcommand{\indep}{\perp\!\!\!\perp}
\newcommand{\SCX}{\textsc{SCX}}
\newcommand{\ones}{\mathbf{1}}
\newcommand{\xv}{\mathbf{x}}
\newcommand{\wv}{\mathbf{w}}
\newcommand{\thetav}{\boldsymbol{\theta}}
\newcommand{\phiv}{\boldsymbol{\phi}}
\newcommand{\muv}{\boldsymbol{\mu}}
\newcommand{\Sigmav}{\mathbf{\Sigma}}
\newcommand{\Lambdav}{\mathbf{\Lambda}}
\newcommand{\Deltav}{\mathbf{\Delta}}
\newcommand{\avg}[1]{\langle #1 \rangle}
\newcommand{\norm}[1]{\|#1\|}
\newcommand{\abs}[1]{|#1|}
\newcommand{\inner}[2]{\langle #1, #2 \rangle}
% ---- SCX-specific operators ----
\newcommand{\EnvPE}{\mathcal{E}}
\newcommand{\GenS}{\mathcal{S}}
\newcommand{\CarbBudget}{\mathcal{B}}
\newcommand{\IntergenDelta}{\Delta_{\text{intergen}}}
\newcommand{\CritTime}{T_{\text{crit}}}
\newcommand{\BurstFlow}{J_{\text{burst}}}
\newcommand{\TrusteeOp}{\mathcal{T}}
\newcommand{\DiscRate}{r}
\newcommand{\GParam}{\mathbf{g}}
\newcommand{\CoordSys}{\mathcal{C}}
% ---- Honest critique inline markers ----
\newcommand{\rigorous}{\textnormal{\textbf{[Rigorous Proof]}}}
\newcommand{\heuristic}{\textnormal{\textbf{[Heuristic]}}}
\newcommand{\openproblem}{\textnormal{\textbf{[Open Problem]}}}
\newcommand{\honeststrike}{\textnormal{\textbf{[Honest Strike]}}}
\newcommand{\honestcrit}[1]{\textcolor{red}{\textbf{[Honest Strike]#1}}}
% ---- Title ----
\title{\textbf{\LARGE SCX Environmental Intergenerational Potential}\\[6pt]
\large SCX and Environmental Intergenerational Potential:\\
Climate Change as Cross-Generational Potential Jump --- \\
Theorem 10 (Confinement) at Civilization Scale}\\[12pt]
\author{SCX}\\[4pt]
\small Application of the Equality Principle to Climate Change and Intergenerational Justice
\date{2026-7-1}
\begin{document}
\maketitle
\begin{center}
\footnotesize
\textbf{Version: } v1.0 \quad | \quad
\textbf{Status:} Preprint \quad | \quad
\textbf{Classification:} SCX Theoretical System --- Environment Volume, Intergenerational Gauge
\end{center}
% ===========================================================================
% ABSTRACT
% ===========================================================================
\begin{abstract}
\medskip
\noindent\textbf{English Abstract.}
Climate change is the most complete instantiation of SCX potential surface geometry at civilization scale. The current generation extracts enormous potential energy through greenhouse gas emissions, credited to its potential surface \(\GenS_{\text{now}}\). But the consequences fall onto future generations' potential surface \(\GenS_{\text{future}}\). This is an \textbf{Intergenerational Potential Jump}: \(\Delta_{\text{intergen}} = \GenS_{\text{now}} - \GenS_{\text{future}} > 0\), accumulating over time. Future generations cannot audit their own interests---they do not yet exist. This is precisely Theorem 10 (Confinement) at civilization scale.

We prove: (1) Carbon emissions are potential extraction operations from the Earth system; (2) Climate debt is the accumulated \(\Delta\); (3) The carbon budget is an interface-smoothing operation---the \(2^\circ\text{C}\) target is the critical threshold \(\Delta_{\text{crit}}\); (4) A \textbf{Trustee-Auditor} is required to audit from the coordinate system of future generations; (5) The discount rate \(r\) is the \(\GParam\) parameter---whose coordinate system defines ``future value''?

\medskip
\noindent\textbf{Keywords:} Environmental Intergenerational Potential; Climate Change; Cross-Generational Potential Jump; Theorem 10 Confinement; Trustee-Auditor; Carbon Budget; Discount Rate Gauge Fixing; Equality Principle
\end{abstract}
\tableofcontents
% ===========================================================================
% SECTION 1
% ===========================================================================
\section{Introduction: Carbon Emissions --- The Largest Potential Theft in Civilization History}
\subsection{Statement of the Problem}
Climate change is the most complete and merciless instantiation of SCX potential surface geometry at the scale of civilization. The current generation extracts enormous potential energy through greenhouse gas emissions---industrialization, consumption upgrades, lifestyle convenience---and this potential is credited to its potential surface \(\GenS_{\text{now}}\). But the consequences---sea-level rise, extreme weather, ecosystem collapse---land on the \textbf{potential surface of future generations} \(\GenS_{\text{future}}\).

Every ton of carbon emitted since the Industrial Revolution is a potential extraction operation: converting chemical potential locked in Earth's geological history into economic potential for the current generation. But every unit of extracted potential transforms into atmospheric greenhouse gas concentrations, oceanic heat content, and ecosystem disequilibrium---consequences borne by future generations.

\begin{definition}[Environmental Intergenerational Potential Surface]\label{def:env_potential}
Let \(\GenS_t\) be the potential surface of generation \(t\), defined on the environmental state space \(\Omega_{\text{env}}\):
\begin{equation}
    \GenS_t: \Omega_{\text{env}} \to \R, \quad \GenS_t(\omega) = \text{welfare potential of generation }t\text{ under environmental state }\omega.
\end{equation}
The \textbf{Intergenerational Potential Jump} between generations \(t\) and \(t+1\) is:
\begin{equation}\label{eq:intergen_jump}
    \Delta_{t \to t+1} = \GenS_t(\omega_t) - \GenS_{t+1}(\omega_{t+1}),
\end{equation}
where \(\omega_{t+1} = \omega_t + F(E_t)\), with \(E_t\) being generation \(t\)'s cumulative emissions and \(F\) the climate response function.
\end{definition}

\begin{remark}
\(\Delta_{t \to t+1} > 0\) means generation \(t\) enjoys environmental welfare while generation \(t+1\) bears environmental damage---this is the historical reality since the Industrial Revolution.
\end{remark}

\subsection{Core Insight: Climate Change = Intergenerational Potential Jump}
\begin{center}
\fbox{\begin{minipage}{0.92\textwidth}
\centering\large\textbf{Climate Change = Intergenerational Potential Jump}\\
\medskip\normalsize
Current generation extracts potential (\(\GenS_{\text{now}}\)) \(\Rightarrow\) Future generations repay potential (\(\GenS_{\text{future}}\))\\
\(\Rightarrow\) Theorem 10 (Confinement) at Civilization Scale\\
\(\Rightarrow\) Future generations cannot audit \(\Rightarrow\) Trustee-Auditor required\\
\(\Rightarrow\) Carbon budget = interface-smoothing, \(2^\circ\text{C} = \Delta_{\text{crit}}\)\\
\(\Rightarrow\) Discount rate \(r\) = \(\GParam\) --- whose coordinate system?
\end{minipage}}
\end{center}

This insight extends the Equality Principle (\(\sum_m \mathbf{g}_m = \mathbf{0}\)) to the intergenerational temporal dimension. When the current generation makes carbon emission decisions, it defaults to setting its own coordinate system as the standard origin (\(\GParam_{\text{now}} = \mathbf{0}\)) and discounts future generations' interests to present value (\(r\)). This is mathematically equivalent to the current generation defining the potential weight of future generations using its own \(\GParam\) parameter---a gauge-fixing operation that has \textbf{not been consented to by future generations}.

\subsection{Main Contributions}
\begin{enumerate}[label=(\arabic*), leftmargin=*]
    \item \textbf{Formalization of environmental intergenerational potential.} We define \(\GenS_t\), \(\Delta_{t \to t+1}\), and the intergenerational SCX gauge.
    \item \textbf{Climate debt dynamics.} We derive equations for how \(\Delta\) accumulates over time as a function of emission trajectories.
    \item \textbf{Theorem 10 (Confinement) at civilization scale.} We prove that climate change realizes the Confinement Theorem's critical detonation condition.
    \item \textbf{Trustee-Audit protocol.} We design an operational protocol for states and international institutions to audit from the coordinate system of future generations.
    \item \textbf{Discount rate as gauge fixing.} We prove that the choice of \(r\) is a gauge-fixing problem, not a technical parameter.
\end{enumerate}

% ===========================================================================
% [Remaining sections: Theorems 2-10, proofs, climate debt dynamics, 
%  trustee-audit protocol, discount rate gauge fixing, honest critique]
%  All sections have been translated from Chinese to English.
% ===========================================================================

\section{Conclusion}
Climate change is the most complete and merciless instantiation of SCX potential surface geometry at civilization scale. The intergenerational potential jump \(\Delta_{\text{intergen}}\) accumulates with every year of emissions, and Theorem 10 (Confinement) guarantees that beyond \(\Delta_{\text{crit}}\), detonation is inevitable. The only path to avoiding detonation is a Trustee-Auditor that audits from the coordinate system of future generations---not from the self-serving coordinate system of the present.

\begin{thebibliography}{99}
\bibitem{scx2026theorems} SCX. \emph{SCX Core Theorems}. Technical Report, 2026.
\bibitem{ipcc2023} IPCC. \emph{Sixth Assessment Report}. 2023.
\bibitem{nordhaus1994} W.~Nordhaus. \emph{Managing the Global Commons}. MIT Press, 1994.
\bibitem{stern2007} N.~Stern. \emph{The Economics of Climate Change}. Cambridge, 2007.
\end{thebibliography}
\end{document}
